\chapter{Vento a 10 m e raffica}
\label{cap:raffica}

\section{Il campo \texorpdfstring{\campo{10fg}}{10fg}}
\label{sec:10fg}

Il campo \campo{10fg} (paramId 49) è la velocità massima del vento a \SI{10}{m} che la simulazione
ha prodotto nell'ora precedente l'istante di validità. Viene da \campo{WSPD10MAX}, uno dei massimi
orari che WRF scrive con \campo{nwp_diagnostics=1}. È un \emph{vento risolto}: non contiene un
termine turbolento né una stima di ciò che la griglia non rappresenta \parencite{issue50}.

La raffica di ERA5 è un'altra grandezza. L'Organizzazione Meteorologica Mondiale definisce la
raffica come il massimo del vento mediato su \SI{3}{s}, un intervallo più breve di un passo temporale
di un modello a \SI{31}{km}, e ERA5 ne pubblica una stima, la \emph{raffica parametrizzata}
\parencite{paramdb}. Nell'IFS la raffica parametrizzata è il vento medio a \SI{10}{m} più un termine
turbolento proporzionale alla velocità d'attrito, più, in condizioni convettive, un termine
proporzionale allo shear del vento nello strato \SIrange{850}{950}{\hPa} \parencite[sez.~3.10.4,
eq.~3.99]{ifs41r2p4}. L'IFS la calcola a ogni passo temporale e ne archivia il massimo dall'ultima
scrittura.

La simulazione regionale riceve da ERA5 lo stato dell'atmosfera e della superficie, non i campi
diagnostici. La tabella standard dei campi ECMWF di WPS 4.1 (\campo{Vtable.ECMWF}) non contiene né il
paramId 49 né il 228029 \parencite{wps41}; la tabella usata da questa simulazione non è stata
verificata. Fra le 200 variabili di un file wrfout nessun campo di raffica si aggiunge a
\campo{WSPD10MAX}, e il campo pubblicato come \campo{10fg} è il più vicino a una raffica che la
simulazione offra.

\section{Confronto con ERA5}
\label{sec:10fg-era5}

La \cref{tab:10fg} riporta il confronto su 16 istanti di validità. \campo{10fg} vale 0,625 volte il
\campo{10fg} di ERA5 e 0,644 volte la raffica parametrizzata istantanea \campo{i10fg}, con
correlazioni di 0,813 e 0,820. Rispetto al vento istantaneo a \SI{10}{m} dello stesso file vale fra
1,004 e 1,032 volte, e in nessuno dei \num{2220273} punti è più basso. Sotto il paramId 49
l'archivio pubblica la velocità del vento a \SI{10}{m}, mentre la raffica parametrizzata di ERA5 vale
circa 1,6 volte il suo vento medio. Il rapporto con il \campo{10fg} di ERA5 resta fra 0,593 e 0,673 in
tutti i 34 istanti del campione di verifica, in ogni mese e a ogni ora, ed è una
differenza di definizione.

\begin{table}[htbp]
  \centering
  \small
  \caption{\campo{10fg} dell'archivio rispetto a tre riferimenti, su 16 istanti di validità.}
  \label{tab:10fg}
  \begin{tabular}{@{}lll@{}}
    \toprule
    Riferimento & CHAPTER / riferimento & Correlazione \\
    \midrule
    \campo{10fg} di ERA5 (raffica parametrizzata, massimo orario) & 0,625 & 0,813 \\
    \campo{i10fg} di ERA5 (raffica parametrizzata, istantanea) & 0,644 & 0,820 \\
    vento istantaneo a \SI{10}{m} di CHAPTER, stesso file & 1,004--1,032 & \\
    \bottomrule
  \end{tabular}
\end{table}

A \SI{3}{km} una parte di ciò che ERA5 deve parametrizzare è risolta, come mostra la velocità
d'attrito sulla terra, fra 1,62 e 2,06 volte quella di ERA5 (\cref{sec:rugosita}). Letto come
raffica, \campo{10fg} sottostima di circa il \SI{38}{\percent}, e nessuna correzione del bias lo rende
una raffica, perché è una grandezza diversa. Con una correlazione di 0,813 con il \campo{10fg} di
ERA5 il campo si presta a ordinare luoghi e ore per ventosità e a descrivere la struttura spaziale
del vento, e fa da limite inferiore della raffica; non si presta a stimare il superamento di una
soglia di raffica.

\section{La finestra oraria}
\label{sec:10fg-finestra}

I messaggi di \campo{10fg} dichiarano un massimo sull'ora precedente (\cref{sec:tempo}). Anche la
raffica parametrizzata oraria di ERA5 è un massimo sull'ora precedente, e la lunghezza della
finestra non è la differenza fra i due campi. La codifica è corretta. WRF azzera \campo{WSPD10MAX}
alla scrittura di ogni file wrfout, ma su sei giorni interi, il 15 luglio 2019 e il 1º e 15 febbraio,
20 marzo, 15 luglio e 15 agosto 2024, per 138 passaggi da un'ora alla successiva, l'azzeramento è
mancato 17 volte, il \SI{12.3}{\percent} dei casi \parencite{issue49}. In quelle ore il campo
contiene il massimo su due ore o più, sotto un messaggio che ne dichiara una.

L'effetto si vede nell'eccesso medio rispetto al vento istantaneo a \SI{10}{m}:
\SIrange{0.068}{0.311}{m/s} nelle ore azzerate, \SIrange{0.499}{0.784}{m/s} in quelle non
azzerate. Su una media di dominio di circa \SI{4.5}{m/s} è una sovrastima di
\SIrange{10}{17}{\percent}. Nelle ore non azzerate il campo si avvicina alla raffica parametrizzata di
ERA5 perché la finestra è più lunga, e il valore non permette di riconoscerle. Neppure gli altri campi
del file lo permettono: i sette massimi orari di WRF si azzerano con lo stesso evento e, quando
l'azzeramento manca, restano tutti insieme, per cui nessuno di essi segnala il mancato azzeramento
degli altri. Riconoscerle richiede il file wrfout dell'ora precedente, che al momento della
conversione non è disponibile. Il meccanismo del mancato azzeramento non è stato stabilito.

Nelle statistiche su una stagione il \SI{12.3}{\percent} delle ore porta una finestra più lunga di
quella dichiarata; per un'ora singola la finestra dichiarata di \campo{10fg} non è garantita.
